\chapter{Học Thêm Và Học Nhóm}
\chapterintro{Danh nghĩa là học, thực tế nhiều bữa là xã giao có đề cương}{Nhóm học chỉ hiệu quả khi người ta sợ bài hơn mê tám chuyện.}

\begin{lucbatbox}{Lục bát: Lập nhóm rất nhanh}
Lập nhóm học tập rất mau
Vừa ngồi đã hỏi trà sữa nơi nào
Sách vở đặt giữa ra sao
Mà câu chuyện chính nhào sang chuyện ngoài
Ba giờ tổng kết rất vui
Bài chưa xong ảnh trong máy đầy hơn
\end{lucbatbox}

\begin{tutuyetbox}{Thất ngôn tứ tuyệt: Học thêm cho yên dạ}
Em đi học thêm đều mỗi chiều
Nghe xong lòng thấy đỡ lo nhiều
Về nhà hỏi lại bài ban nãy
Tinh thần trắng tựa một trang phiếu
\end{tutuyetbox}

\begin{ghichu}
Học nhóm thất bại không vì đông người, mà vì phần thích nói thường lấn phần thật sự muốn hiểu.
\end{ghichu}

\ngancach

\begin{lucbatbox}{Lục bát: Chia việc rất đều}
Bạn kia chụp ảnh đề cương
Bạn này giữ nhịp chuyện thường rất hay
Bạn kia ngồi bấm điện thoại
Riêng người làm thật im lìm cuối bàn
Nộp xong ai cũng hân hoan
Nhóm mình cố gắng vô vàn trên môi
\end{lucbatbox}

\begin{tutuyetbox}{Thất ngôn tứ tuyệt: Học nhóm bằng đồ ăn}
Người mang bánh tráng với trà chanh
Sách vở ngồi im ở cạnh mình
No bụng là điều ai cũng thấy
Còn bài nhìn lại vẫn mông mênh
\end{tutuyetbox}